
\chapter{\\INTRODUCTION}
\section{Background}

Over the past decade, Vietnamese higher education has undergone a profound structural shift driven by national digital transformation policy. Vietnam's National Education Development Strategy 2021--2030, enacted under Decision No.~1373/QD-TTg, establishes concrete targets for the integration of digital technology into teaching and learning at all levels of the education system~\cite{vietnam_strategy2021}. This strategy is not merely a technological upgrade but a fundamental rethinking of how educational content is created, distributed, and consumed across the nation. Institutions affiliated with Vietnam National University -- Ho Chi Minh City (VNU-HCM) are at the forefront of this effort, tasked with pioneering new methodologies that can scale effectively to meet the growing demands of a rapidly developing knowledge economy. The VNU-HCM MOOC platform, publicly available at \url{https://courses.vnuhcm.edu.vn}, has delivered thousands of course modules to learners across the country and serves as the primary digital learning infrastructure for the university~\cite{vnuhcm_mooc}. 

Despite this significant progress in establishing the foundational infrastructure, the gap between infrastructure availability and effective pedagogical use remains wide. Surveys conducted at HCMIU indicate that while 89\% of instructors express interest in creating interactive digital content, only 23\% report confidence in using current authoring tools. Faculty members cite interface complexity, the steep learning curve associated with plugin-based workflows, and a lack of dedicated instructional design support as the primary barriers to adoption. This gap is not unique to Vietnam; comparative research on e-learning quality in developing university contexts consistently identifies technical usability as a decisive factor in adoption rates~\cite{hadullo2017}. The challenge is therefore not a lack of will or institutional infrastructure, but a lack of purpose-built tools that meet educators where they are, seamlessly bridging the gap between pedagogical intent and technological execution.

Interactive video has emerged as one of the most effective media for asynchronous and hybrid instruction because it combines the explanatory power of video with the engagement mechanisms of formative assessment~\cite{zhang2006}. Rather than passively consuming a recorded lecture—a mode of learning often criticized for inducing cognitive passivity and high attrition rates—students interacting with embedded questions must retrieve and apply knowledge in real time. This active retrieval mechanism is well-supported by the evidence base for active learning~\cite{freeman2014}. Merkt et al.\ demonstrated that videos with interactive control features---particularly stop-and-answer moments---produced significantly better knowledge retention and transfer scores than equivalent linear content when learners were actively engaged rather than passively watching~\cite{merkt2011}. Vural similarly found that embedding questions at strategic timestamps within educational videos improved post-test scores, enhanced long-term memory retention, and drastically reduced learner dropout in online environments~\cite{vural2013}.

The pedagogical design of these embedded questions is best guided by Bloom's revised taxonomy, which provides a principled framework for targeting specific cognitive levels from basic recall (\emph{remember}) through synthesis and evaluation (\emph{create})~\cite{anderson2001}. When questions are distributed systematically across these cognitive levels and positioned precisely at moments where a key concept has just been introduced, the extraneous cognitive load imposed by the interruption is minimised, while germane cognitive load---the load associated with meaningful schema construction and integration into long-term memory---is maximised~\cite{sweller1988,chandler1991,paas2003}.

H5P (HTML5 Package) is an open-source framework for creating, sharing, and embedding richly interactive content in web environments~\cite{h5p2023,lumieducation2023}. Originally developed by Joubel AS in 2013, the H5P ecosystem has expanded to encompass over 50 content types, ranging from foundational assessments like Multiple Choice and True/False questions to sophisticated branching scenarios, virtual tours, and interactive timelines. A core design principle of the framework is that all content is entirely self-contained: each H5P package exists as a single .h5p ZIP archive containing its own JavaScript libraries, styling assets, media files, and a structured JSON configuration representing the content parameters. This inherent portability allows H5P packages to be trivially imported into any H5P-aware Learning Management System (LMS), such as Moodle~\cite{moodle2023}, Canvas, and Blackboard, via standard web integration protocols or the LTI~1.3 standard~\cite{lti2019}.

H5P Interactive Video stands as the most widely utilized content type within higher education. The framework wraps a video source—whether a direct upload or a streaming URL like YouTube—in an intelligent H5P player that pauses playback at pre-defined timestamps to present interactive overlays. These elements require explicit learner input before playback can resume, effectively transforming a passive viewing experience into an active dialogue. The specific interaction types supported and highly relevant to this implementation include Multiple Choice, True/False, Fill in the Blanks, Drag Text, Mark the Words, and Matching. Each of these types maps to a dedicated H5P JavaScript library with a strictly defined content parameter schema, facilitating the programmatic generation of H5P content from structured JSON representations.

Despite H5P's profound capabilities, the dominant deployment model in Vietnamese universities relies heavily on the legacy WordPress H5P plugin. This plugin was initially designed for small-scale personal blogs or individual websites rather than large-scale institutional content production systems. Creating a single interactive video through this pathway requires an instructor to execute a convoluted, multi-step process: (1)~log in to a WordPress administration panel; (2)~navigate complex menus to create or locate a ``Post''; (3)~activate the H5P shortcode block; (4)~navigate to a completely separate H5P content editor iframe; (5)~upload or link the video source; (6)~manually locate each desired timestamp through a built-in player interface; (7)~configure an interaction type by navigating through several nested, dense settings screens; and (8)~save the configuration and return to the main post. This workflow spans multiple disjointed interfaces, involves web publishing concepts (posts, shortcodes, plugin taxonomy) that lie entirely outside an educator's pedagogical domain, and takes an average of 45 to 60 minutes for a first-time user to complete even a rudimentary three-question interactive video.

Research into educational technology adoption identifies the Technology Acceptance Model (TAM) as a highly useful lens for understanding why educators accept or reject tools~\cite{davis1989}. Perceived ease of use and perceived usefulness together explain a substantial fraction of the variance in adoption intention. When a tool imposes high extraneous cognitive load on every interaction---as the WordPress H5P workflow demonstrably does through its fragmented interface---perceived ease of use collapses, and systemic adoption fails to materialize regardless of the tool's underlying educational potential or institutional backing~\cite{sweller1988,chandler1991}.

The recent, rapid maturation of large language models (LLMs) has opened an entirely new frontier for educational content authoring: the ability to analyze a video's transcript and generate contextually appropriate, pedagogically structured interactive questions automatically~\cite{brown2020,anthropic2024}. If an instructor can transition from the role of a manual content creator to that of a content curator—reviewing and accepting high-quality, AI-generated suggestions rather than composing each question and configuring its metadata from scratch—the cognitive overhead of interactive video creation falls dramatically. The same AI pipeline that extracts a transcript, classifies its segments by topic and complexity, and returns a structured list of H5P interactions can also apply Bloom's taxonomy labels to each suggested question. This automated pedagogical tagging helps instructors achieve a deliberate spread of cognitive difficulty throughout the video, ensuring comprehensive assessment without requiring specialist training in instructional design.

This thesis explores whether combining a purpose-built, user-centred web authoring interface with an advanced AI pipeline—powered by Anthropic Claude for cloud-based inference and a locally hosted Ollama inference server for privacy-preserving, cost-effective generation—can produce a measurable, transformative improvement in educator productivity and satisfaction when creating H5P interactive videos.

% ─────────────────────────────────────────────────────────────────────────────
\section{Problem Statement}

The central problem this thesis addresses is the stark mismatch between the proven educational potential of H5P interactive video and the practical usability barriers of existing authoring environments for university instructors.

The problem manifests through four interacting challenges that compound to limit adoption:

\begin{enumerate}
  \item \textbf{Workflow complexity and fragmentation.} Current WordPress-based H5P workflows require educators to navigate at least six distinct screens and understand web publishing concepts that are completely irrelevant to their core pedagogical task. Task analysis of the existing workflow measured an astonishing 47 interface interactions (clicks, scrolls, form entries) across six screens to produce just one interactive video. Furthermore, these interfaces suffer from 3.2-second median page-load times, severely interrupting the instructor's creative flow and causing significant frustration.
  
  \item \textbf{High extraneous cognitive load.} The mental effort required to operate the technical interface reduces the cognitive resources available for actual pedagogical design~\cite{chandler1991,mayer2009}. Instructors report that managing the tool's intricate settings, timestamp alignments, and shortcode placements consumes more attention than designing the questions themselves. This inversion of effort fundamentally breaks the authoring experience.
  
  \item \textbf{Low institutional adoption rates.} As a direct consequence of the overwhelming workflow complexity and cognitive load, adoption of H5P interactive video at VNU-HCM has been artificially limited to a small subset of highly technically proficient faculty. Senior instructors, as well as those in humanities and social science disciplines who lack extensive IT backgrounds, are systematically excluded from utilizing a tool that empirical research consistently demonstrates to be highly educationally effective~\cite{davis1989,hadullo2017}.
  
  \item \textbf{Absence of intelligent AI assistance.} No existing, freely available H5P authoring tool provides integrated, native AI assistance for question generation. Instructors who could benefit the most from AI support---particularly those with limited formal instructional design experience---have no streamlined mechanism for obtaining context-aware question suggestions tied directly to the transcript and content of their specific video.
\end{enumerate}

The \textbf{primary research question} driving this thesis is therefore:

\begin{quote}
  \itshape How can a purpose-built web platform, intelligently combining rigorous user-centred interface design with an advanced, multi-modal AI-driven transcript analysis pipeline, improve the efficiency, pedagogical quality, and institutional adoption of H5P interactive video authoring by university instructors?
\end{quote}

Four supporting research questions are established to systematically guide the investigation:

\begin{enumerate}[label=RQ\arabic*.]
  \item What are the specific usability barriers and friction points in current H5P authoring workflows that most strongly reduce task completion rates and user satisfaction?
  \item How can foundational user-centred design principles be systematically applied to reduce extraneous cognitive load in timestamp-based interactive question authoring?
  \item What full-stack technical architecture best supports scalable, real-time AI-assisted question generation and seamless H5P package assembly within a unified web environment?
  \item To what extent does the newly developed, purpose-built platform quantitatively and qualitatively outperform the legacy WordPress-based baseline on measurable usability, efficiency, and productivity metrics?
\end{enumerate}

% ─────────────────────────────────────────────────────────────────────────────
\section{Scope and Objectives}

The thesis pursues five concrete, measurable objectives in order to comprehensively answer the research questions outlined above.

\begin{enumerate}
  \item \textbf{Platform Development and Architecture.} Design, engineer, and deploy a robust full-stack web application that enables instructors to seamlessly upload or link a video. The platform must provide an intuitive visual timeline interface to add H5P interactive elements at precise timestamps, allow for real-time content previewing, and ultimately export a fully standards-compliant \texttt{.h5p} package suitable for direct, flawless import into the VNU-HCM MOOC infrastructure and other LTI-compliant Learning Management Systems.

  \item \textbf{Advanced AI Pipeline Integration.} Architect and integrate a dual-provider, fault-tolerant AI pipeline. This pipeline must accurately extract a timestamped transcript from the uploaded video (leveraging YouTube automatic captions, uploaded WebVTT/SRT subtitle files, or an integrated Whisper audio transcription service). The system must securely transmit the transcript to a large language model (supporting both the cloud-based Anthropic Claude API and a locally hosted Ollama instance for sensitive content), and return a strictly validated JSON payload containing H5P interaction suggestions. These suggestions must include precise timestamps, appropriate content types, distractor options, and a pedagogical Bloom's taxonomy classification.

  \item \textbf{Rigorous User-Centred Design.} Apply standard user-centred design (UCD) methodology throughout the entire interface design and iterative prototyping process. All design decisions must be empirically grounded in identified instructor needs, mapped to detailed user personas, and continuously validated against established Nielsen usability heuristics.

  \item \textbf{Comprehensive Security Implementation.} Implement a robust security baseline rigorously aligned to the OWASP Top Ten vulnerabilities. This includes secure JWT-based authentication, strict role-based access control (RBAC) authorisation, comprehensive server-side input validation using Zod schemas, enforced transport layer security (HTTPS), and aggressive rate limiting to prevent API abuse and denial-of-service vectors.

  \item \textbf{Empirical Usability Evaluation.} Design and conduct a controlled, comparative usability study. This study must statistically measure critical metrics including task completion rate, time-on-task, error frequency, and System Usability Scale (SUS) scores against the established WordPress H5P baseline. The evaluation will also capture qualitative data through think-aloud protocols and post-task interviews to provide deep contextual insights into the user experience.
\end{enumerate}

% ─────────────────────────────────────────────────────────────────────────────
The project scope is explicitly bounded by the following constraints to ensure feasibility and focus.

\textbf{Content type scope.} The platform specifically targets the six H5P interaction types most pedagogically relevant to interactive video assessment: Multiple Choice, True/False, Fill in the Blanks, Drag Text, Mark the Words, and Matching. While all six are fully supported for manual authoring through the interface, the AI pipeline is restricted to generating Multiple Choice, True/False, Fill in the Blanks, and Matching suggestions automatically. These specific types possess highly deterministic schema structures that can be reliably and strictly validated via Zod, ensuring the AI output never breaks the H5P player.

\textbf{Video source scope.} The platform natively accepts direct video file uploads in standard web formats (MP4, WebM, and MOV) up to a hard limit of 500~MB, accommodating standard lecture lengths. It also seamlessly accepts YouTube video URLs via the H5P YouTube integration. For YouTube-sourced videos, transcripts are extracted instantly from the platform's automatic or manual captions API. For uploaded video files, instructors may manually provide a WebVTT caption file; alternatively, the system can automatically invoke the \texttt{faster-whisper} machine learning model locally to perform highly accurate audio-to-text transcription.

\textbf{User scope.} The primary target user profile is a university instructor or teaching assistant at VNU-HCM. System administration, configuration, and monitoring are handled by a distinct secondary user role (Admin), which possesses elevated privileges to manage user accounts, review system usage analytics, and export security audit logs.

\textbf{Integration scope.} The platform is designed to export standard, universally compatible H5P packages that can be imported into any LTI-compatible LMS globally. While an LTI~1.3 launch route is stubbed in the backend architecture for direct LMS integration and deep-linking, complex student-level gradebook passback and granular learning-analytics reporting via xAPI are explicitly excluded from this iteration and are designated as future work.

The following limitations are acknowledged within the study design:

\begin{itemize}
  \item \textbf{H5P content types.} The platform deliberately does not attempt to support the massive, full catalogue of 50+ diverse H5P content types. Complex, multi-layered types such as Branching Scenarios, Course Presentations, and Virtual Tours are strictly out of scope for both manual authoring and AI generation.
  \item \textbf{Evaluation sample limitations.} Usability testing was conducted utilizing ten senior Information Technology students serving as educator proxies, rather than actual university faculty members. While this proxy approach is a standard, cost-effective, and widely accepted practice for early-stage formative usability evaluation~\cite{nielsen1994}, it inherently introduces a limitation regarding ecological validity and generalizability to older, less technically fluent demographics.
  \item \textbf{Mobile authoring constraints.} The platform is engineered as a responsive web application optimized specifically for desktop and large tablet browsers, reflecting the reality of complex video editing tasks. While mobile authoring is technically possible due to responsive CSS, it was not a primary design priority and was not formally evaluated in the usability study.
  \item \textbf{Concurrent user load.} The system architecture was stress-tested using load simulation tools at up to 50 concurrent active users. Massive-scale institutional deployments serving hundreds or thousands of simultaneous video editors across a distributed university system were not evaluated and would likely require a more robust Kubernetes-based microservices architecture.
  \item \textbf{AI transcription quality dependency.} The pedagogical quality of AI-generated question suggestions depends entirely and directly on the accuracy of the underlying video transcript. Poor audio quality, heavy background noise, or strong regional accents in uploaded videos will significantly degrade Whisper transcription accuracy, which will consequently hallucinate or degrade the relevance of the downstream AI question generation.
\end{itemize}

% ─────────────────────────────────────────────────────────────────────────────
\section{Assumption and Solution}

This thesis makes several substantial contributions to both the academic research landscape of educational technology and practical software engineering application.

\begin{enumerate}
  \item \textbf{A reproducible, scalable platform architecture.} The thesis meticulously documents a modern, full-stack architecture combining React~18, Zustand state management, Express.js, SQLite/Sequelize ORM, the highly complex Lumieducation H5P server library, and a resilient dual-provider AI pipeline. This detailed architectural blueprint enables exact reproduction, modification, and large-scale institutional adoption by other universities.
  \item \textbf{A novel AI-assisted authoring workflow.} The sophisticated data pipeline that seamlessly transitions a raw video transcript through structured, context-aware LLM prompting, into a strictly Zod-validated, Bloom's-taxonomy-labelled list of H5P interactive suggestions constitutes a significant and novel integration of generative AI into the established H5P open-source ecosystem.
  \item \textbf{Advanced Prompt Engineering for structured H5P JSON output.} The thesis details the highly optimized system prompts, few-shot examples, and JSON validation strategies developed specifically to constrain and coerce models like Claude and Ollama into returning complex JSON payloads that perfectly conform to strict H5P content type specifications. This provides a reusable, robust design pattern for generating structured AI output in other educational content management systems.
  \item \textbf{Empirical usability evidence.} By conducting a rigorous comparative evaluation against the widely used WordPress H5P baseline, the thesis provides strong quantitative evidence that a purpose-built, user-centred platform can substantially reduce authoring time, significantly lower error rates, and vastly improve overall task completion and satisfaction. This contributes heavily to the growing evidence base advocating for user-centred design in educational technology development.
  \item \textbf{Open-source contribution.} The complete platform codebase—including all intricate AI service routing logic, prompt template definitions, database schemas, and H5P integration middleware—is made openly available on public repositories. This commitment to open science supports further academic research, collaborative improvement, and free institutional deployment in resource-constrained environments.
\end{enumerate}

The practical significance and real-world impact of this work is therefore profound and threefold. First, by drastically lowering the technical authoring barrier and cognitive load, the platform directly enables a far broader range of VNU-HCM faculty—including those with minimal IT skills—to independently produce highly engaging interactive content. This democratisation of content creation directly improves the quality of online learning experienced by tens of thousands of students. Second, the successful AI-assisted workflow serves as a powerful proof-of-concept, demonstrating that the often theoretical gap between cutting-edge generative AI research and practical, reliable educational tool development can be successfully bridged with intelligent software engineering. This encourages further institutional investment in similar AI integrations. Third, the platform aligns perfectly with Vietnam's overarching national digital transformation goals~\cite{vietnam_strategy2021} by providing a sovereign, locally deployable, open-source solution that drastically reduces reliance on expensive, proprietary commercial educational software.

% ─────────────────────────────────────────────────────────────────────────────
\section{Structure of thesis}

The remainder of this comprehensive report is meticulously organised as follows to guide the reader through the research, engineering, and evaluation phases.

\textbf{Chapter~2 -- Literature Review} surveys the deep academic and technical foundations underpinning the work. It covers the evolution of e-learning in Vietnamese higher education, provides a highly detailed analysis of the H5P framework and its architectural limitations, and reviews empirical research on interactive video pedagogy. The chapter also explores cognitive load theory, user-centred design principles, the rapid evolution of large language models, advanced prompt engineering techniques, cloud-native web architecture paradigms, and LMS integration standards. It concludes with a critical analysis of related systems, clearly identifying the specific research gaps this project intends to address.

\textbf{Chapter~3 -- Methodology} describes the overarching research design, detailing the rigorous requirements analysis process, the user-centred design (UCD) methodology applied during interface conceptualization and wireframing, the agile software development approach utilized for engineering, and the strict empirical evaluation protocol designed to accurately measure usability outcomes.

\textbf{Chapter~4 -- System Design and Implementation} presents the complete, in-depth system architecture. It details the relational database design, complex backend API routing implementation, dynamic frontend React component structure, and state management strategies. Crucially, it provides a deep technical dive into the AI pipeline's mechanics, the intricacies of the H5P server integration, and the comprehensive security implementation safeguarding the platform.

\textbf{Chapter~5 -- Results and Evaluation} rigorously reports the outcomes of the comparative usability study. It presents quantitative task metrics, statistical analysis of SUS scores, and a thematic analysis of qualitative participant feedback. Furthermore, it reports on the technical performance of the AI pipeline, analyzing latency, generation accuracy, and overall system responsiveness under load.

\textbf{Chapter~6 -- Conclusion and Future Work} synthesises all research findings, critically assesses the project's contributions in direct relation to the initial research questions, transparently discusses the methodological and technical limitations of the work, and proposes an extensive, actionable roadmap for future research, pedagogical enhancement, and platform scalability.
